% 本文件由 main.tex 输入，作为“模型设定与变量说明”章节，不单独编译。

\section{模型设定与变量说明}

\subsection{数据来源与样本说明}

本文以地级市—年份为研究单元，样本期为2015—2024年。选择地级市层面展开分析，主要是因为政府产业引导基金通常以注册地、出资地和服务地方产业为重要运行边界，地方债务压力和财政约束也主要在地方政府层面体现。与企业层面数据相比，地级市面板更适合识别地方财政债务约束下政府产业引导基金与区域创新产出之间的关系。

本文使用的创新产出数据来自 CNRDS 地方专利数据，并整理为地级市—年份层面的发明申请量、实用新型申请量和专利申请总量。政府产业引导基金基础信息来自清科基金目录和投中基金目录，本文将两类目录合并去重后，依据基金注册地手工匹配到地级市层面，进而构造各城市基金当年设立数量、累计设立数量和累计设立规模。基金出资结构主要来自投中基金目录中的 LP 和出资人信息，用以识别政府资本、国资平台、政府产业引导基金、FOF 以及其他社会资本出资。基金投资事件信息主要来自清科投资事件数据，本文据此识别政府引导基金的投资时间、投资阶段、投资次数和投资金额，并按基金注册地汇总到地级市—年份层面。地方债务数据依据各地方统计局和统计年鉴披露信息手工整理，控制变量主要来自 CEIC、CNKI 等城市统计资料。

样本处理遵循现有回归面板的统一口径。回归前按城市和年份进行去重，模型中加入城市固定效应和年份固定效应，并按城市维度聚类稳健标准误。由于部分债务变量、基金投资事件变量和机制变量存在缺失，不同回归表中的样本量会随变量口径变化，本文在实证结果表中分别报告各规格观测值数量。

\subsection{变量定义}

本文被解释变量为城市创新产出。正文主结果使用专利申请量而非授权量，原因在于申请量能够及时地反映城市创新活动变化。具体而言，本文分别使用发明申请量、实用新型申请量和专利申请总量衡量创新产出，其中发明申请量更接近高质量创新口径，专利申请总量则作为综合创新产出口径。涉及对数形式时，统一使用 $\ln(1+\text{专利量})$ 处理。

核心解释变量为政府产业引导基金规模。基准回归同时考察基金当年设立数量、累计设立数量和累计设立规模。相较于基金数量口径，累计设立规模更直接反映某地截至当年通过引导基金形成的政策性资本供给总量，也更贴近本文关注的政府资本支持强度和资源配置能力，因此本文将其作为正文主线变量，并将基金当年设立数量和累计设立数量作为替代基金口径。

调节变量为地方债务压力。本文采用债券余额合计与公共财政收入之比衡量地方债务压力，该值越高，表示地方偿债压力和财政约束越强。相较于财政收支压力等流量指标，该指标更直接刻画地方政府存量债务负担及其对财政空间、风险承担能力的约束。为缓解同期反向因果和共同冲击的影响，正文主规格优先采用滞后一期债务压力；当期债务压力作为补充口径。财政收支压力更多反映年度预算平衡紧张程度，因此在本文中仅作补充讨论，不作为正文主调节变量。

机制变量围绕三条理论路径设置。第一，耐心资本属性用早期投资事件占比衡量。若投资阶段为种子期或初创期，则记为早期投资，早期投资事件占比等于早期投资事件数除以该市引导基金当年投资总次数。相较于早期投资金额占比，事件占比更直接反映基金是否将投资机会配置到早期项目，也较少受少数大额投资事件影响。第二，社会资本撬动效率用社会资本出资与政府资本出资之比及其变化衡量，其中社会资本出资剔除国资平台、政府产业引导基金和 FOF 后识别。第三，融资约束缓解效应使用地区金融发展水平作为反向代理，正文主口径为存款/GDP，并采用反双曲正弦变换，因为该指标更适合刻画地区金融资源供给能力的存量水平；金融发展增速口径用于补充反映短期改善速度。金融发展水平越高，表示地区金融供给能力越强、融资约束相对越弱，但该指标不能直接等同于企业层面的 SA、KZ 或 WW 融资约束指数。

\begin{table}[htbp]
  \centering
  \caption*{主要变量定义与构造}
  {\fangsong\zihao{-5}
  \begin{tabularx}{0.96\linewidth}{l l X}
    \toprule
    变量类型 & 变量名称 & 构造说明 \\
    \midrule
    创新产出 & 发明申请量 & CNRDS 地方专利数据中地级市—年份发明专利申请数量，用于衡量高质量创新产出。 \\
    创新产出 & 实用新型申请量 & CNRDS 地方专利数据中地级市—年份实用新型专利申请数量，用作补充创新口径。 \\
    创新产出 & 专利申请总量 & 发明、实用新型和外观设计等专利申请数量汇总，用于衡量综合创新产出。 \\
    基金变量 & 基金当年设立数量 & 清科和投中基金目录合并去重后，按注册地和设立年份汇总到城市层面。 \\
    基金变量 & 基金累计设立数量 & 截至当年各地级市累计设立的政府产业引导基金数量。 \\
    基金变量 & 基金累计设立规模 & 截至当年各地级市累计设立的政府产业引导基金规模，正文主解释变量。 \\
    债务压力 & 债务压力 & 债券余额合计除以公共财政收入，数值越高表示地方债务约束越强。 \\
    债务压力 & 滞后一期债务压力 & 债务压力的一期滞后项，正文主调节变量。 \\
    耐心资本 & 早期投资事件占比 & 种子期和初创期投资事件数占该市引导基金当年投资事件总数的比例。 \\
    社会资本 & 社会资本撬动效率 & 非政府社会资本出资与政府资本出资之比，用于衡量政府资本放大效应。 \\
    金融发展 & 存款/GDP & 金融机构存款与地区生产总值之比，用作地区融资约束的反向代理。 \\
    \bottomrule
  \end{tabularx}

  \par\smallskip
  {\songti\zihao{6}注：涉及对数变量时，规模变量统一采用 $\ln(x+1)$ 处理；金融发展主机制采用反双曲正弦变换。}
  }
\end{table}

表1报告主要变量的描述性统计。样本中城市专利申请总量均值为13109.508件，发明申请量均值为4683.297件。基金累计设立规模均值为4.439亿元，但标准差达到16.662亿元，说明政府产业引导基金在城市之间分布差异较大。债务压力均值为233.176，滞后一期债务压力均值为199.112，也显示地方债务约束存在明显地区差异。早期投资事件占比均值为0.340，表明政府引导基金投资事件中约三分之一投向种子期或初创期项目。

\begin{table}[htbp]
  \centering
  \caption{主要变量描述性统计}
  {\fangsong\zihao{-5}
  \resizebox{\linewidth}{!}{%
  \begin{tabular}{l r r r r r}
    \toprule
    变量 & N & 均值 & 标准差 & 最小值 & 最大值 \\
    \midrule
    发明申请量 & 2720 & 4683.297 & 13016.643 & 6.000 & 195388.000 \\
    实用新型申请量 & 2720 & 6420.434 & 13622.132 & 21.000 & 146592.000 \\
    专利申请总量 & 2720 & 13109.508 & 29764.731 & 34.000 & 323014.000 \\
    基金当年设立数量 & 2770 & 2.803 & 7.215 & 0.000 & 99.000 \\
    基金累计设立数量 & 2770 & 14.266 & 39.521 & 0.000 & 458.000 \\
    基金累计设立规模（亿元） & 2629 & 4.439 & 16.662 & 0.000 & 279.266 \\
    债务压力 & 1904 & 233.176 & 484.577 & 0.746 & 4654.410 \\
    滞后一期债务压力 & 1898 & 199.112 & 422.855 & 0.806 & 4129.294 \\
    早期投资事件占比 & 957 & 0.340 & 0.338 & 0.000 & 1.000 \\
    社会资本撬动效率 & 986 & 0.789 & 2.015 & 0.000 & 34.050 \\
    金融发展水平 & 2500 & 1.638 & 0.772 & 0.000 & 20.100 \\
    经济发展水平 & 2098 & 5.453 & 0.864 & 3.081 & 8.590 \\
    财政科技支出 & 1971 & 6.343 & 1.418 & 2.757 & 11.060 \\
    人口规模 & 2450 & 5.980 & 0.748 & 3.252 & 11.975 \\
    第二产业规模 & 2098 & 4.554 & 0.886 & 2.272 & 7.238 \\
    外资利用水平 & 1681 & 10.184 & 2.082 & 1.386 & 14.776 \\
    \bottomrule
  \end{tabular}%
  }

  \par\smallskip
  {\songti\zihao{6}注：基金累计设立规模按亿元报告。经济发展水平、财政科技支出、人口规模、第二产业规模和外资利用水平为 $\ln(x+1)$ 口径。社会资本撬动效率报告有限值观测，剔除因政府出资为0形成的非有限值。}
  }
\end{table}

控制变量沿用现有回归设定，主要包括经济发展水平、财政科技支出、人口规模、产业结构和外资利用水平。经济发展水平用地区生产总值对数衡量，财政科技支出用于控制地方财政对科技创新的直接支持强度，常住人口对数用于控制人口规模和市场规模，第二产业增加值对数用于控制城市产业基础，实际利用外资额对数用于控制对外开放程度。上述控制变量均按 $\ln(x+1)$ 处理，以缓解规模变量右偏分布并保留取值为零的观测。

\subsection{模型设定}

为检验政府产业引导基金是否促进城市创新，本文首先估计如下双向固定效应模型：

\begin{equation}
Innovation_{it}=\alpha+\beta Fund_{it}+\gamma Controls_{it}+\mu_i+\lambda_t+\varepsilon_{it}.
\end{equation}
\ERJWhere{$Innovation_{it}$ 表示城市 $i$ 在年份 $t$ 的创新产出，$Fund_{it}$ 表示政府产业引导基金变量，$Controls_{it}$ 为城市层面控制变量，$\mu_i$ 和 $\lambda_t$ 分别表示城市固定效应和年份固定效应。若 $\beta$ 显著为正，则说明政府产业引导基金扩张与城市创新产出提升相一致。}

在此基础上，本文引入债务压力与基金规模的交互项，考察地方债务压力是否削弱基金的创新促进效应：

\begin{equation}
Innovation_{it}=\alpha+\beta_1 Fund_{it}+\beta_2 Debt_{i,t-1}
+\beta_3(Fund_{it}\times Debt_{i,t-1})+\gamma Controls_{it}+\mu_i+\lambda_t+\varepsilon_{it}.
\end{equation}
\ERJWhere{$Debt_{i,t-1}$ 表示滞后一期债务压力，$Fund_{it}\times Debt_{i,t-1}$ 为核心交互项。本文关注 $\beta_3$ 的符号和显著性。若 $\beta_3<0$，说明债务压力越高，同等规模政府产业引导基金对创新产出的边际促进作用越弱。}

机制检验围绕耐心资本属性、社会资本撬动效率和融资约束缓解效应展开。耐心资本和金融发展机制主要采用“债务压力削弱基金影响机制变量，机制变量进一步影响创新产出”的路径检验：

\begin{equation}
M_{it}=\alpha+\theta_1(Fund_{it}\times Debt_{i,t-1})+\gamma Controls_{it}+\mu_i+\lambda_t+\nu_{it},
\end{equation}

\begin{equation}
Innovation_{i,t+k}=\alpha+\rho_1(Fund_{it}\times Debt_{i,t-1})+\rho_2 M_{it}
+\gamma Controls_{it}+\mu_i+\lambda_t+\eta_{it}.
\end{equation}
\ERJWhere{$M_{it}$ 表示机制变量，$k$ 表示创新产出的滞后实现期。耐心资本机制中 $M_{it}$ 为早期投资事件占比，理论预期为 $\theta_1<0$ 且 $\rho_2>0$；金融发展机制中 $M_{it}$ 为金融发展水平，理论预期同样为 $\theta_1<0$ 且 $\rho_2>0$。}

社会资本撬动效率机制不写作普通中介，而写作调节式机制或承接机制。相应模型在结果方程中进一步加入基金规模与社会资本撬动效率的交互项：

\begin{equation}
Innovation_{it}=\alpha+\delta_1(Fund_{it}\times Debt_{i,t-1})
+\delta_2(Fund_{it}\times SocCap_{it})+\gamma Controls_{it}+\mu_i+\lambda_t+\xi_{it}.
\end{equation}
\ERJWhere{$SocCap_{it}$ 表示社会资本撬动效率或其变化。若债务压力降低撬动效率，同时 $\delta_2>0$，且加入机制项后 $\delta_1$ 的绝对值下降，则说明债务压力对基金创新效应的抑制部分通过社会资本撬动效率这一承接渠道体现。}

以上模型均采用城市固定效应和年份固定效应，并按城市维度聚类稳健标准误。本文的识别重点不是宣称完全解决所有内生性问题，而是在固定城市不随时间变化的特征、年份共同冲击以及主要城市层面控制变量后，比较同一城市在不同年份基金规模、债务压力和创新产出的共同变化关系。后续滞后债务压力和滞后创新响应的检验，用于进一步缓解反向因果和同期冲击的担忧。
